\BeleskeID{09_2-s007}
% element: 09_2-s007:e05
Kako ulazni napon raste, izlazni napon opada. Opterećenje ostaje u zasićenju, a pogonski tranzistor ulazi u omsku oblast. Ravnoteža struja pre zamene napona glasi
\[\frac{k_{n,L}}2(V_{GS,L}-V_{Tn,L})^2=\frac{k_{n,D}}2\left(2V_{DS,D}(V_{GS,D}-V_{Tn,D})-V_{DS,D}^2\right)\]
U oblasti gornjeg ulaznog praga diferencira se ravnoteža struja i zamenjuje nagib $-1$:
\[0=-k_{n,D}\left(2(V_I-V_{Tn,D})-V_O\right)+k_{n,D}V_O(2+1)\]
Otuda se dobija druga jednačina $V_I=V_{Tn,D}+2V_O$. Zajedničkim rešavanjem sa ravnotežom struja dobijaju se $V_{O(IH)}$ i $V_{IH}$. Oba rezultata pripadaju istom paru radnih tačaka.
